%%%%%%%% ICML 2026 EXAMPLE LATEX SUBMISSION FILE %%%%%%%%%%%%%%%%%

\documentclass{article}

% Recommended, but optional, packages for figures and better typesetting:
\usepackage{microtype}
\usepackage{graphicx}
\usepackage{booktabs}

% hyperref makes hyperlinks in the resulting PDF.
% If your build breaks (sometimes temporarily if a hyperlink spans a page)
% please comment out the following usepackage line and replace
% \usepackage{icml2026} with \usepackage[nohyperref]{icml2026} above.
\usepackage{hyperref}

% Attempt to make hyperref and algorithmic work together better:
\newcommand{\theHalgorithm}{\arabic{algorithm}}

% Authored preprint. Use \usepackage{icml2026} for a blind submission.
\usepackage[preprint]{icml2026}
\newcommand{\icmlpreprinttext}{} 
\renewcommand{\icmlpreprinttext}{}

% If accepted, instead use the following line for the camera-ready submission:
% \usepackage[accepted]{icml2026}

\usepackage{amsmath}
\usepackage{amssymb}
\usepackage{mathtools}
\usepackage{amsthm}


% if you use cleveref..
\usepackage[capitalize,noabbrev]{cleveref}

%%%%%%%%%%%%%%%%%%%%%%%%%%%%%%%%
% THEOREMS
%%%%%%%%%%%%%%%%%%%%%%%%%%%%%%%%
\theoremstyle{plain}
\newtheorem{theorem}{Theorem}[section]
\newtheorem{proposition}[theorem]{Proposition}
\newtheorem{lemma}[theorem]{Lemma}
\newtheorem{corollary}[theorem]{Corollary}
\theoremstyle{definition}
\newtheorem{definition}[theorem]{Definition}
\newtheorem{assumption}[theorem]{Assumption}
\theoremstyle{remark}
\newtheorem{remark}[theorem]{Remark}

% The \icmltitle you define below is probably too long as a header.
% Therefore, a short form for the running title is supplied here:
\icmltitlerunning{Remoteness-Aware Control of Negative Off-Policy Updates}
\setlength{\abovedisplayskip}{3pt}
\setlength{\belowdisplayskip}{3pt}

\begin{document}

\twocolumn[
  \icmltitle{Breaking the Curse of Repulsion: Remoteness-Aware Control of Negative Off-Policy Updates}

  % It is OKAY to include author information, even for blind submissions: the
  % style file will automatically remove it for you unless you've provided
  % the [accepted] option to the icml2026 package.

  % List of affiliations: The first argument should be a (short) identifier you
  % will use later to specify author affiliations Academic affiliations
  % should list Department, University, City, Region, Country Industry
  % affiliations should list Company, City, Region, Country

  % You can specify symbols, otherwise they are numbered in order. Ideally, you
  % should not use this facility. Affiliations will be numbered in order of
  % appearance and this is the preferred way.
  \icmlsetsymbol{equal}{*}

  \begin{icmlauthorlist}
    \icmlauthor{Jie Jiang}{equal,comp}
    \icmlauthor{Yusen Huo}{equal,comp}
    \icmlauthor{Changping Wang}{comp}
    \icmlauthor{Jun Zhang}{comp}
  \end{icmlauthorlist}
  \icmlaffiliation{comp}{Tencent Inc, China}

  \icmlcorrespondingauthor{Jun Zhang}{neoxzhang@tencent.com}

  % You may provide any keywords that you find helpful for describing your
  % paper; these are used to populate the "keywords" metadata in the PDF but
  % will not be shown in the document
  \icmlkeywords{Off-Policy Reinforcement Learning, Negative Feedback, Policy Optimization, Stability}

  \vskip 0.3in
]

% this must go after the closing bracket ] following \twocolumn[ ...

% This command actually creates the footnote in the first column listing the
% affiliations and the copyright notice. The command takes one argument, which
% is text to display at the start of the footnote. The \icmlEqualContribution
% command is standard text for equal contribution. Remove it (just {}) if you
% do not need this facility.

% Use ONE of the following lines. DO NOT remove the command.
% If you have no special notice, KEEP empty braces:
% \printAffiliationsAndNotice{}  % no special notice (required even if empty)
% Or, if applicable, use the standard equal contribution text:

\printAffiliationsAndNotice{\icmlEqualContribution}


\begin{abstract}
Off-policy policy optimization reuses historical behavior, including
negative-advantage samples that suppress known failures. We show that repeated
reuse can turn this useful signal into excessive repulsion: as the learner
moves away from a historical negative action, subsequent updates make that
action increasingly remote without necessarily reducing its update strength.
Our aggregate theory characterizes the resulting transition from a stable
displacement beyond the positive-only target to persistent drift and the loss
of finite stable equilibria; controlled strength sweeps show that an
intermediate displacement can improve held-out reward. The relevant
learner-relative coordinate is squared standardized distance for Gaussian
policies and surprisal for categorical policies. We introduce Dynamic
Remoteness-Aware Policy Optimization (DRPO), which leaves the negative update
unchanged in the near field and exponentially attenuates its remote tail.
DRPO restores eventual inward Gaussian drift for every fixed finite
negative-to-positive mass ratio, yields an explicit ultimate-bound radius, and
changes categorical support suppression from exponential to polynomial
probability decay. External diagnostics and controlled interventions isolate
remoteness-dependent policy geometry as a source of negative-update
amplification and show that selective tapering can remove its far-field effect
without discarding useful local feedback.
\end{abstract}

\section{Introduction}
\label{sec:introduction}

Policy-based methods learn from signed feedback: positive updates attract the
policy toward actions estimated to be preferable, whereas negative updates
suppress actions estimated to be undesirable. Negative feedback can therefore
encode useful boundaries and exclusion information that positive-only learning
lacks \citep{zhu2025negative,song2025good}. In off-policy settings, however,
the same historical sample may be reused after earlier updates have already
changed its probability under the current policy. This raises a single
question: when does repeated reuse turn useful negative feedback into
excessive repulsion?

We identify learner-relative remoteness as the state variable that links one
reuse step to the next. A negative update lowers the probability of its
historical action, increasing the action's remoteness and thereby changing the
next policy-score response. Informally, we call historical actions below the
control threshold the \emph{near field} and those beyond it the \emph{far
field}; Section~\ref{sec:drpo} makes this distinction operational through
\(D\leq\tau\) and \(D>\tau\). The policy family determines the resulting tail
law: Gaussian mean scores grow with standardized distance, whereas categorical
logit scores remain bounded even as selected-action probabilities approach the
simplex boundary. Under isolated reuse, uncontrolled Gaussian remoteness grows
geometrically and categorical action probability decays exponentially; an
exponential taper reduces the corresponding far-field growth to logarithmic
order. This is the curse of repulsion: complying with a negative update can
make its next reuse more remote without making the reused signal self-limiting.

What remains unresolved is whether remoteness itself amplifies negative
updates independently of sample quality. Existing methods control off-policy
instability through filtering, importance weighting, probability-aware gates,
and tapered updates
\citep{novati2019remember,fakoor2020p3o,peng2019advantage,
schulman2017proximal,fujimoto2019off,kumar2020conservative,
arnal2025asymmetric,leroux2025tapered}, but remoteness is usually confounded
with advantage magnitude and data quality. We therefore match those factors
while varying only remoteness, then intervene separately on near- and far-field
negative updates to test their downstream effects.

Our aggregate analysis places these per-sample effects in competition with
positive attraction. When positive mass dominates, repulsion can displace the
policy beyond the positive-only target while preserving a finite stable
equilibrium. At balance, the restoring force disappears and persistent drift
emerges; under negative dominance, no finite stable equilibrium remains.
These regimes characterize policy geometry and stability, not task utility,
which must be determined empirically.

Motivated by this transition, we introduce Dynamic Remoteness-Aware Policy
Optimization (DRPO). DRPO recomputes remoteness under the current policy,
leaves the negative branch unchanged in the near field, and applies an
exponential taper beyond the threshold. The taper dominates every finite-order
far-field response, restores eventual inward Gaussian drift for any fixed
finite negative-to-positive mass ratio, and makes categorical suppression
self-limiting. Section~\ref{subsec:distance_dependent_utility} states the important
boundary of this design: remoteness alone cannot identify remote samples that
remain aligned with true task improvement.

The experiments separate existence, identification, transmission, and control.
External diagnostics first test whether the predicted pattern appears in
realistic neural policies; controlled environments then isolate its source,
trace its causal effect on policy behavior, and compare selective tapering
against positive-only, global-scaling, and polynomial-tail alternatives.

Our contributions are summarized as follows:
\begin{itemize}
    \item We characterize how repeated negative reuse changes
    policy-relative remoteness and derive aggregate regimes for stable
    extrapolation, persistent drift, and the loss of finite stable
    equilibria.

    \item We derive a control-order hierarchy and a selective exponential
    taper with Gaussian ultimate-boundedness and categorical
    self-limiting-suppression guarantees.

    \item We use matched external diagnostics and controlled interventions
    to separate coefficient magnitude from policy geometry, identify a
    far-field causal pathway, and test the stability--utility trade-off of
    selective tapering.
\end{itemize}


% ------------------------------------------------------------------
% SECTION 2: PRELIMINARIES
% ------------------------------------------------------------------
\section{Related Work}
\label{sec:related_work}

\textbf{Learning from Negative or Suboptimal Behavior.}
A broad line of policy-learning methods distinguishes behavior according
to its estimated quality. Advantage-weighted regression, AWAC, implicit
Q-learning, critic-regularized regression, and behavior-proximal
objectives fit actors more strongly toward actions estimated to be
valuable under a learned critic
\citep{peng2019advantage,nair2020awac,kostrikov2021offline,
wang2020critic,zhuang2023behavior}. Positive-filtered or filtering-based
variants represent a conservative endpoint, where selected negative or
low-quality updates are removed rather than reweighted. At the same
time, failed or suboptimal behavior can remain informative: it can
suppress competing bad modes, sharpen local decision boundaries, or
improve data efficiency
\citep{novati2019remember,fakoor2020p3o,arnal2025asymmetric}. These
findings support the view that negative feedback is not merely noise, but
also carries information whose usefulness depends on how it is used.

\textbf{Off-Policy, Stale, and Learner-Relative Updates.}
This issue becomes especially important under off-policy reuse, where
historical behavior remains fixed while the learner continues to change.
Off-policy methods commonly control behavior--learner mismatch through
importance weighting, clipping, trust regions, replay mechanisms, and
behavior regularization
\citep{schulman2017proximal,haarnoja2018soft,levine2020offline,
fujimoto2019off,kumar2020conservative}. Recent asymmetric or tapered
off-policy objectives further make updates depend on current-policy
probability or surprisal
\citep{arnal2025asymmetric,leroux2025tapered}. These works establish
that mismatch and rarity matter; our distinction concerns how
learner-relative remoteness is created and amplified by repeated
negative reuse. This actor-side view is complementary to conservative
offline RL, which primarily addresses value extrapolation and
policy-support mismatch.

\textbf{Tapered Updates and Optimistic Reweighting.}
TOPR uses an asymmetric tapered importance ratio on negative off-policy
updates and provides its own stability analysis
\citep{leroux2025tapered}. Optimistic or best-case distributional
optimization also has a formal precedent in likelihood approximation
and policy optimization
\citep{nguyen2019optimistic,song2020optimistic}. Our construction is
related but distinct: it reweights a finite measure of negative-update
mass according to learner-relative remoteness and may reduce that
measure's total mass. It is therefore an analogue of optimistic
divergence-ball reweighting, not a standard ambiguity set over normalized
probability distributions.
% ------------------------------------------------------------------
% SECTION 2: PRELIMINARIES
% ------------------------------------------------------------------

\section{Problem Setup}
\label{sec:problem_setup}

% ------------------------------------------------------------------
% SECTION 2: PRELIMINARIES
% ------------------------------------------------------------------

\subsection{Off-Policy Signed Actor Updates}

We consider policy optimization from a historical update distribution
$\nu$ over state--action pairs, which may represent an offline dataset,
a replay buffer, or trajectories collected by earlier or stale behavior
policies. Let $\pi_\theta(a \mid s)$ denote the current policy and let
$\widehat{A}(s,a)$ denote the advantage-like signal used by the actor
update. During an actor step, the update distribution and the associated
advantage estimates are treated as fixed with respect to $\theta$. The resulting empirical actor field is
\begin{equation}
    \mathbf{F}(\theta)
    =
    \mathbb{E}_{(s,a)\sim\nu}
    \left[
        \widehat{A}(s,a)
        \nabla_\theta \log \pi_\theta(a \mid s)
    \right],
    \label{eq:empirical_actor_field}
\end{equation}
with the corresponding update
\(\theta_{t+1}=\theta_t+\eta\mathbf{F}(\theta_t)\).
Repeated reuse means that the same historical samples and signed feedback
are evaluated under an evolving current policy;
Section~\ref{sec:theory} studies the actor-side dynamics induced by this
reuse.

To separate favorable and unfavorable feedback, we define
\begin{equation}
\begin{aligned}
    \widehat{A}^{+}(s,a)
    &= \max\{\widehat{A}(s,a),0\}, \\
    \widehat{A}^{-}(s,a)
    &= \max\{-\widehat{A}(s,a),0\}.
\end{aligned}
\end{equation}
The actor field can then be decomposed into positive and negative
components:
\begin{equation}
\begin{aligned}
    \mathbf{F}(\theta)
    ={}&
    \mathbb{E}_{\nu}
    \left[
        \widehat{A}^{+}(s,a)
        \nabla_\theta \log \pi_\theta(a \mid s)
    \right] \\
    &-
    \mathbb{E}_{\nu}
    \left[
        \wi